% Bewertungspaket zum Gesamttest «Quadratische Gleichungen» (GF 2.2b) — Quelle des PDFs
% (python3 scripts/build-lp-pdf.py leitprogramme/quadratische-gleichungen/).
% Ganze Punkte je Teilschritt; Werte und Fehlerwerte mit python3 nachgerechnet (zahlen.py, 11.10.2026).
\documentclass[11pt]{article}
\usepackage{../lp-druck}
\renewcommand{\lpname}{Quadratische Gleichungen}
\renewcommand{\lpbereich}{Grundlagenfach}
\renewcommand{\lpteilgebiet}{2.2b}
\renewcommand{\lpfuss}{Bewertungspaket}
\newcolumntype{P}{>{\centering\arraybackslash}p{10mm}}
\newenvironment{raster}{\par\smallskip\noindent\tabularx{\linewidth}{@{}>{\raggedright\arraybackslash}p{38mm}P X@{}}\toprule
  \small\textsc{Teilschritt} & \small\textsc{P} & \small\textsc{Musterlösung}\\\midrule}{\bottomrule\endtabularx\par}
\newcommand{\fehler}[1]{\par\smallskip{\small\textcolor{lprot}{\bfseries Typische Fehler:} #1}}
\newcommand{\E}{\,{\footnotesize\bfseries\color{lpgruen}(E)}}
\begin{document}
\lpkopf{Bewertungspaket zum Gesamttest}{}

\begin{lpachtung}{Erst lösen, dann öffnen}
Dieses Paket enthält die Musterlösung. Wer es vor dem Lösen liest, misst am Ende nur, wie gut das Abschreiben gelingt.
\end{lpachtung}

\section*{1 · So gehst du vor}
\begin{enumerate}[leftmargin=*,itemsep=2pt]
\item \textbf{Gesamttest lösen} — vollständig, mit Rechenweg, ohne dieses Paket und ohne Taschenrechner.
\item \textbf{Fotografieren oder scannen}, gut lesbar, eine Seite pro Bild. \textbf{Kein Name} auf den Bildern.
  Handyfotos tragen oft unsichtbar Standort, Zeit und Gerät mit (Metadaten). Lade darum lieber einen Scan oder ein
  Bildschirmfoto des Fotos hoch, oder entferne vorher die Standortangaben.
\item \textbf{Bewerten lassen.} Ob du dafür eine KI nutzt und welche, entscheidest du selbst und in eigener Verantwortung —
  je nachdem, was dir aufgrund deines Alters und deiner persönlichen Situation erlaubt ist (Altersgrenzen und
  Nutzungsbedingungen des Anbieters, allenfalls Einverständnis deiner Eltern). Lade nur deine Lösung und dieses PDF hoch,
  keine persönlichen Daten, und füge den Auftrag aus Abschnitt~2 ein. Ohne KI kannst du dich mit Abschnitt~3 selbst bewerten.
\item \textbf{Rückmeldung prüfen.} Stimmt, was die KI aus deiner Handschrift gelesen hat? Eine KI kann sich irren —
  im Zweifel gilt das Raster in Abschnitt~3.
\item \textbf{Gezielt wiederholen} nach Abschnitt~4.
\end{enumerate}

\section*{2 · Auftrag an die KI}
\begin{tcolorbox}[colback=lplinie!25,colframe=lplinie,boxrule=0.5pt,arc=1mm,fontupper=\small]
Du bist Korrektorin oder Korrektor für Mathematik an einer Schweizer Berufsmaturitätsschule. Im Anhang findest du (1)~das Bewertungspaket zum Gesamttest «Quadratische Gleichungen» mit Musterlösung und Punkteraster und (2)~Fotos meiner handschriftlichen Lösung.\par\medskip
Geh so vor:\par
1. Schreib zu jeder Aufgabe G1 bis G7 zuerst ab, was du in meiner Lösung liest — Rechenweg und Ergebnis. Was du nicht sicher lesen kannst, markierst du mit [unsicher]. Nicht raten.\par
2. Vergib die Punkte genau nach dem Raster im Bewertungspaket (Abschnitt 3), ganze Punkte je Zeile. Ziehe einen Fehler nur einmal ab: Rechne ich mit einem falschen Zwischenergebnis richtig weiter, gibt es die Folgepunkte. Ausnahme: Zeilen mit \textbf{(E)} gibt es nur für das richtige Ergebnis, nie als Folgepunkt. Die Zeilen «Typische Fehler» sagen, welche Rasterzeile bei diesem Fehler 0 Punkte gibt — sie sind kein zusätzlicher Abzug.\par
3. Andere richtige Lösungswege zählen voll, ausser die Aufgabe verlangt ein bestimmtes Verfahren (G3 quadratische Ergänzung, G4 Mitternachtsformel, G2(b) Faktorisieren — dort regelt das Raster unter «Typische Fehler» den anderen Weg). Gleichwertige Schreibweisen zählen voll: Dezimalpunkt oder Dezimalkomma, Bruch oder Dezimalzahl ($\tfrac32 = 1.5$), gekürzt oder ungekürzt, $\mathbb{L} = \{-2;\, 8\}$ oder «$x = -2$ oder $x = 8$» oder $x_{1,2} = \ldots$, die Elemente in beliebiger Reihenfolge, $\mathbb{L} = \{\,\}$ oder «keine Lösung». Zeilen mit \textbf{(E)} sind Ergebnispunkte: Das Ergebnis muss stimmen \emph{und} aus einem erkennbaren Weg hervorgehen. Zeilen mit «Folgezeile» geben den Punkt auch für richtiges Weiterrechnen mit einem eigenen falschen Zwischenwert.\par
4. Begründe jeden Abzug in einem Satz und nenne die Fehlerart (zum Beispiel «Vorzeichen von $-4ac$» oder «Fall $k = -2$ übersehen»). Schreib die Musterlösung nicht ab — zeig mir, wo mein Fehler liegt.\par
5. Gib zum Schluss eine Tabelle «Aufgabe | Punkte | Begründung», die Gesamtpunktzahl (von 25) und — nach der Selbsteinschätzung in Abschnitt 4 — welches Kapitel des Leitprogramms ich wiederholen soll. Keine Note.\par
6. Fehlt ein Foto oder ist eine Aufgabe unlesbar, sag es, statt sie zu bewerten.
\end{tcolorbox}

\needspace{16\baselineskip}
\section*{3 · Musterlösung und Punkteraster}
\textbf{Allgemein:} Ganze Punkte je Zeile. Ein Fehler wird nur einmal abgezogen (Folgefehler); derselbe Begriffsfehler
nochmals gibt keinen Folgepunkt. In der Spalte P: \textbf{(E)}~= Ergebnispunkt für ein Ergebnis, das direkt aus den Angaben
folgt — nur das richtige Ergebnis zählt, nie als Folgepunkt. Ohne Zeichen = die Zeile bewertet den Weg; steht «Folgezeile»,
zählt auch ein Ergebnis, das aus einem eigenen falschen Zwischenwert richtig weitergerechnet ist. Der ganze Test ist ohne Taschenrechner gestellt; exakte Werte
(Brüche, Wurzeln) sind verlangt, wo die Aufgabe es sagt. «Typische Fehler» nennt, welche Zeile 0 Punkte gibt, und die Punkte, die bleiben.
\textbf{Teilrichtige Zeilen:} Verlangt eine Zeile mehrere Angaben und stimmt nur ein Teil, gibt die Zeile 0 Punkte —
ausser die Zeile nennt ausdrücklich eine Aufteilung. \textbf{Fehlende Lösung:} Fehlt in einer Lösungsmenge eine Lösung (zum
Beispiel nur die positive Wurzel), gibt die Ergebniszeile 0 Punkte — kein Folgepunkt.
\textbf{Komma:} Wer das Dezimalkomma benutzt, muss die Elemente von $\mathbb{L}$ mit Strichpunkt trennen. Steht «$\{1,5\}$»
ohne Strichpunkt und ist aus dem Weg nicht klar, ob $1.5$ oder $1$ und $5$ gemeint ist: [unsicher] melden, nicht raten.

\begin{lpaufgabe}{G1}{Binom und Wurzel}{3 P}
\begin{raster}
(a) Binom & 1 & $x^2 + 10x + 25 = (x + 5)^2$, also $(x + 5)^2 = 9$\\
(a) Lösung & 1 & $x + 5 = \pm 3$: $\mathbb{L} = \{-8;\, -2\}$. Folgezeile: Wer mit einem falschen Binom richtig Wurzel zieht, erhält den Punkt.\\
(b) Begründung & 1 & Es wäre $(x + 5)^2 = -1$; ein Quadrat ist nie negativ, also $\mathbb{L} = \{\,\}$. Gleichwertig: $x^2 + 10x + 26 = 0$ hat $D = 100 - 104 = -4 < 0$.\\
\end{raster}
\fehler{Nur $x + 5 = 3$, also $\{-2\}$: Lösungszeile 0 (fehlende Lösung) $\to$ 2 von 3\,P.
$(x - 5)^2 = 9$ und daraus richtig $\{2;\, 8\}$: Binomzeile 0, Lösungszeile 1 (Folgefehler) $\to$ 2 von 3\,P.
Ausmultipliziert und $x^2 + 10x + 16 = 0$ mit Zweiklammersatz oder Formel gelöst: anderer richtiger Weg, Binom- und Lösungszeile voll.
(b) Nur «keine Lösung» ohne Grund: Zeile 0.}
\end{lpaufgabe}

\begin{lpaufgabe}{G2}{Eine fremde Lösung und eine Zerlegung}{4 P}
\begin{raster}
(a) Begründung & 1 & Verloren geht $x = 4$ (Probe: $0 = 0$): Teilen durch $(x - 4)$ setzt $x \neq 4$ voraus.\\
(a) Lösung & 1 & $(x - 4)^2 - 3(x - 4) = 0$, $(x - 4)(x - 7) = 0$: $\mathbb{L} = \{4;\, 7\}$. Gleichwertig: ausmultipliziert $x^2 - 11x + 28 = 0$ und faktorisiert, oder Fallunterscheidung «$x = 4$ prüfen, dann für $x \neq 4$ teilen».\\
(b) Faktorform & 1 & Erst auf null: $x^2 - x - 30 = 0$, dann $(x - 6)(x + 5) = 0$ (Lösungen mit Produkt $-30$ und Summe $1$).\\
(b) Lösung und Kontrolle & 1 & $\mathbb{L} = \{-5;\, 6\}$ mit Vieta: $-5 + 6 = 1 = -p$, $(-5) \cdot 6 = -30 = q$. Beides nötig; ohne Kontrolle 0. Folgezeile: Eine richtig gerechnete Kontrolle, die einen eigenen Fehler aufdeckt und benennt, gibt den Punkt.\\
\end{raster}
\fehler{(a) Nur «$x = 4$ fehlt» ohne Grund: Begründungszeile 0. Ausmultipliziert, aber falsch geordnet: Lösungszeile 0.
(b) Nullprodukt bei rechts $30$ («$x = 30$ oder $x - 1 = 30$»): beide (b)-Zeilen 0.
(b) Vorzeichen $(x + 6)(x - 5) = 0$, also $\{-6;\, 5\}$: Faktorzeile 0; deckt die Kontrolle den Fehler auf (Summe $-1 \neq 1$) und wird er benannt, Lösungszeile 1.
(b) Mitternachtsformel statt Faktorisieren ($D = 121$): Faktorzeile 0, Lösungszeile 1, wenn die Kontrolle dasteht.}
\end{lpaufgabe}

\begin{lpaufgabe}{G3}{Quadratische Ergänzung}{3 P}
\begin{raster}
Ordnen und Ergänzung & 1 & $x^2 + 6x = -2$; beidseitig $+\left(\tfrac{6}{2}\right)^2 = +9$: $(x + 3)^2 = 7$\\
Wurzel & 1 & $x + 3 = \pm\sqrt{7}$ (Folgezeile)\\
Lösung & 1 & $\mathbb{L} = \{-3 - \sqrt7;\, -3 + \sqrt7\}$ (Folgezeile). Nur gerundet ($\approx -5.65$; $-0.35$) ohne exakte Form: 0 — die Aufgabe verlangt exakt.\\
\end{raster}
\fehler{Nullprodukt bei rechts $-2$ («$x = -2$ oder $x + 6 = -2$», also $-2$ und $-8$ — beide falsch): alle Zeilen 0.
$+9$ nur links addiert, also $(x + 3)^2 = -2$ und «keine Lösung»: erste Zeile 0; Wurzel- und Lösungszeile je 1 als Folgefehler («Quadrat nie negativ» richtig angewandt) $\to$ 2 von 3\,P.
Statt $+9$ eine andere Zahl ergänzt (z.\,B. $+36$) und links trotzdem $(x + 3)^2$ geschrieben: erste Zeile 0; weiter wie Folgefehler. Steht links kein Quadrat, gibt es keine weiteren Punkte.
Nur $x + 3 = \sqrt7$: Wurzelzeile 0, Lösungszeile 0 (fehlende Lösung).
Mitternachtsformel statt Ergänzung ($D = 28$, $x = \tfrac{-6 \pm 2\sqrt7}{2}$): erste und Wurzelzeile 0, Lösungszeile 1, wenn exakt richtig.}
\end{lpaufgabe}

\begin{lpaufgabe}{G4}{Erst ordnen, dann die Formel}{3 P}
\begin{raster}
Ordnen & 1 & $2x^2 - 2x + x - 1 = 2$, also $2x^2 - x - 3 = 0$\\
$D$ & 1 & $a = 2$, $b = -1$, $c = -3$: $D = 1 + 24 = 25$ (Folgezeile)\\
Lösung & 1 & $x = \dfrac{1 \pm 5}{4}$: $\mathbb{L} = \{-1;\, 1.5\}$ (Folgezeile)\\
\end{raster}
\fehler{Nullprodukt bei rechts $2$ («$2x + 1 = 2$ oder $x - 1 = 2$», also $0.5$ und $3$ — beide falsch): alle Zeilen 0.
Nicht auf null gebracht, mit $2x^2 - x - 1 = 0$ weitergerechnet: Ordnungszeile 0; $D = 9$ und $\{-0.5;\, 1\}$ als Folgefehler je 1 $\to$ 2 von 3\,P.
Mittleres Glied vergessen ($2x^2 - 3 = 0$): Ordnungszeile 0; $D = 24$ und $x = \pm\tfrac{\sqrt{24}}{4}$ als Folgefehler je 1.
$-b$ falsch eingesetzt ($\tfrac{-1 \pm 5}{4}$, also $\{-1.5;\, 1\}$): Lösungszeile 0.
Faktorisieren statt Formel ($(2x - 3)(x + 1) = 0$): Die Aufgabe verlangt die Formel — Ordnungszeile 1, $D$-Zeile 0, Lösungszeile 1, wenn richtig.}
\end{lpaufgabe}

\begin{lpaufgabe}{G5}{Verfahren begründen und einen Fehler finden}{5 P}
\begin{raster}
(a)(i) & 1 & Binom erkennen: $x^2 - 16x + 64 = (x - 8)^2$. Gleichwertig: Zweiklammersatz mit den Lösungen $8$ und $8$. Verfahren und ein Grund nötig.\\
(a)(ii) & 1 & Typ: Ausmultipliziert $x^2 + 4x + 4 = x^2 + 8$, die $x^2$ heben sich weg — die Gleichung ist \emph{linear} ($4x = 4$, $x = 1$; das Lösen ist nicht verlangt). Verfahren: umformen wie bei linearen Gleichungen. «Linear» mit diesem Grund nötig.\\
(a)(iii) & 1 & Mitternachtsformel: Faktor $3$ vor $x^2$ und alle drei Glieder da, oder: keine Zerlegung mit ganzen Zahlen ($D = 61$). Verfahren und ein Grund nötig.\\
(b) Fehler & 1 & Bei $c = -21$ ist $-4ac = -4 \cdot 1 \cdot (-21) = +84$: richtig ist $D = 16 + 84 = 100$.\\
(b) Lösung & 1 & $x = \dfrac{4 \pm 10}{2}$: $\mathbb{L} = \{-3;\, 7\}$ (Folgezeile). Gleichwertig: $(x - 7)(x + 3) = 0$.\\
\end{raster}
\fehler{Verfahren ohne Begründung: Zeile 0. Mitternachtsformel oder quadratische Ergänzung bei (i), auch mit «geht immer»: Zeile 0 — gefragt ist das schnellste Verfahren.
(ii) «quadratisch» oder ein quadratisches Verfahren genannt: Zeile 0.
(iii) Zweiklammersatz oder Wurzelziehen: 0. Quadratische Ergänzung nach Teilen durch $3$: 0 — sie ist hier nicht schneller als die Formel.
(b) Nur «$D$ ist falsch» ohne den Vorzeichenfehler zu nennen: Fehlerzeile 0; Lösungszeile trotzdem voll, wenn richtig.}
\end{lpaufgabe}

\begin{lpaufgabe}{G6}{Parameter im linearen und im konstanten Glied}{3 P}
\begin{raster}
$D(k)$ & 1\E & $a = 1$, $b = k$, $c = k$: $D = k^2 - 4k = k\,(k - 4)$\\
Fälle & 1 & $D = 0$ bei $k = 0$ und $k = 4$: eine Lösung. $0 < k < 4$: $D < 0$, keine. $k < 0$ oder $k > 4$: $D > 0$, zwei. Alle drei Fälle nötig (Folgezeile).\\
Lösungen bei $D = 0$ & 1 & $k = 0$: $x^2 = 0$, $x = 0$. $k = 4$: $(x + 2)^2 = 0$, $x = -2$. Beide nötig (Folgezeile).\\
\end{raster}
\fehler{Aus $k^2 - 4k = 0$ nur $k = 4$ (durch $k$ geteilt): Fälle- und Lösungszeile 0 (fehlender Fall) $\to$ 1 von 3\,P.
Fälle vertauscht (zwei Lösungen für $0 < k < 4$): Fälle-Zeile 0.
$D = k^2 - 4$ (das zweite $k$ als $1$ gelesen): $D$-Zeile 0; Fälle mit $k = \pm 2$ richtig weitergedacht: Fälle-Zeile 1 als Folgefehler; Lösungszeile 0 — bei $k = \pm 2$ hat die gegebene Gleichung keine doppelte Lösung ($x^2 + 2x + 2$ hat $D = -4$), das Einsetzen zeigt es.}
\end{lpaufgabe}

\begin{lpaufgabe}{G7}{Parameter vor dem quadratischen Glied}{4 P}
\begin{raster}
Fall $k = 0$ & 1 & Faktor vor $x^2$ null: $4x = 0$, linear, genau eine Lösung $x = 0$\\
$D(k)$ & 1\E & $k \neq 0$: $D = 4^2 - 4 \cdot k \cdot k = 16 - 4k^2$\\
Fälle & 1 & $D = 0$ bei $k = 2$ und $k = -2$: eine Lösung. $-2 < k < 2$ und $k \neq 0$: zwei. $k < -2$ oder $k > 2$: keine. Alle drei nötig, beim Fall «zwei» samt Ausschluss $k \neq 0$ (Folgezeile).\\
Lösungen bei $D = 0$ & 1 & $k = 2$: $2(x + 1)^2 = 0$, $x = -1$. $k = -2$: $-2(x - 1)^2 = 0$, $x = 1$. Beide nötig (Folgezeile).\\
\end{raster}
\fehler{Fall $k = 0$ übersehen, nur mit $D$ gerechnet: erste Zeile 0; Fälle-Zeile 0, weil der Ausschluss $k \neq 0$ fehlt $\to$ höchstens 2 von 4\,P.
Bei $k = 0$ mit $D(0) = 16 > 0$ «zwei Lösungen»: erste Zeile 0.
Nur $k = 2$ (aus $k^2 = 4$): Fälle- und Lösungszeile 0 (fehlender Wert).
$D = 16 - 4k$ ($c$ als $1$ statt $k$): $D$-Zeile 0; Fälle mit $k = 4$ richtig weitergedacht: Fälle-Zeile 1 als Folgefehler; Lösungszeile 0, weil der zweite Wert fehlt.}
\end{lpaufgabe}


\section*{4 · Selbsteinschätzung}
\begin{tabularx}{\linewidth}{@{}l X@{}}\toprule
22 – 25 P & Die geprüften Teile sitzen. Wo du Punkte verloren hast: das Kapitel dieser Aufgabe nochmals (Zuordnung unten).\\
17 – 21 P & Das Kapitel, in dem du die meisten Punkte verloren hast, nochmals: Tüfteln und Übungen.\\
11 – 16 P & Zurück zu den Kapiteln aller Aufgaben, in denen du Punkte verloren hast.\\
0 – 10 P & Zurück zu Kapitel 1 und von dort der Reihe nach weiter.\\\midrule
\multicolumn{2}{@{}p{\linewidth}@{}}{\small\textbf{Aufgabe $\to$ Kapitel:} G1 $\to$ Kapitel 1 und 2 (Wurzel ziehen, Binom); G2 $\to$ Kapitel 2 und 4 (Faktorisieren, erst auf null); G3 $\to$ Kapitel 3 und 4 (quadratische Ergänzung, ordnen); G4 $\to$ Kapitel 3 und 4 (Formel, ordnen); G5 $\to$ Kapitel 4, 1 und 3 (Typ und Verfahren, Diskriminante); G6, G7 $\to$ Kapitel 5 (Parameterdiskussion)}\\\bottomrule
\end{tabularx}
\end{document}
